
@@M0@@
De plus l'additivité de @@M1@@ entraine que @@M2@@, et les relations du corollaire sont vérifiées.

Si @@M3@@ est la @@M4@@-algèbre engendrée par des générateurs @@M5@@, soumis aux relations~:
@@M6@@
il existe donc un homomorphisme surjectif de @@M7@@ sur @@M8@@. D'après 4.4, cet homomorphisme transforme une base de @@M9@@ en une base de @@M10@@, c'est donc un isomorphisme, ce qui achève la démonstration du corollaire.

Traitons le cas général d'un fibré en drapeaux de type @@M11@@. Le fibré en drapeaux usuel @@M12@@ peut s'obtenir à partir de @@M13@@ en construisant successivement les fibrés en drapeaux des @@M14@@ (4.2). Il en résulte que @@M15@@, considéré comme module sur @@M16@@, admet une base formée des éléments considérés dans 4.3.i), les @@M17@@ étant les classes des faisceaux analogues aux @@M18@@ intervenant dans le scindage des images inverses des @@M19@@ sur @@M20@@ (notations de 4.5 et 4.7). Donc @@M21@@ est un sous-anneau de @@M22@@ ; de plus, l'additivité de @@M23@@ entraine que~:
@@M24@@
et @@M25@@ contient la sous-@@M26@@-algèbre de @@M27@@ engendrée par les @@M28@@. Soit @@M29@@ cette sous-algèbre ; d'après 4.3.i), et ce qui vient d'être dit, on a dans @@M30@@ une famille d'éléments qui forme une base à la fois pour la structure de @@M31@@-module et pour la structure de @@M32@@-module ; par suite, @@M33@@, ce qui montre que @@M34@@ est engendré par les @@M35@@. L'idéal des relations entre les @@M36@@ est alors donné par 4.3.ii), compte
